\section{Introduction}
\label{sec:introduction}

\subsection{The Alpha Hyperparameter in Multi-scale MCMC}

Multi-scale MCMC (\ms{}) is a Search-layer strategy for redistricting
chains introduced in G.11~\citep{dellaLibera2026multiscale}.
The algorithm alternates between two Markov kernels: a \emph{fine-level}
\recom{} move on the full census-tract graph, and a \emph{coarse-level}
move on the county-aggregated graph, projected back to the tract level via
a rebalancing pass.
Coarse moves allow the chain to escape fine-scale local optima by proposing
large-scale boundary rearrangements at once; fine moves maintain detailed
feasibility.

The key hyperparameter is $\alpha \in (0,1)$: the probability of making a
coarse move at each step.
Setting $\alpha$ too low underuses the coarse level, leaving the chain
effectively equivalent to standard \recom{}.
Setting $\alpha$ too high causes the chain to repeatedly propose coarse moves
that fail the rebalancing check, wasting steps.

\subsection{State-Specific Sensitivity}

The optimal $\alpha$ is not universal.
It depends on three structural properties of the state:

\paragraph{County/tract ratio.}
A coarse county graph with $C$ counties has roughly $C$ nodes; a fine
tract graph has $T$ tracts.
The ratio $T/C$ measures how much larger a county-level proposal is compared
to a tract-level proposal.
For Texas ($C = 254$, $T = 5{,}265$, ratio $\approx 20.7$) each coarse move
proposes a very large rearrangement, and a moderate $\alpha \approx 0.30$
gives enough coarse coverage without overwhelming the chain.
For Iowa ($C = 99$, $T = 825$, ratio $\approx 8.3$) the coarse moves are
comparatively small, and the rebalancing constraint is harder to satisfy
with many districts, making a lower $\alpha \approx 0.15$ more effective.
California is an extreme case: Los Angeles County alone holds approximately
25\% of the state's population, making county-level proposals for $k = 52$
districts very difficult to rebalance, so an even lower $\alpha$ is preferred.

\paragraph{District count $k$.}
Higher $k$ means the rebalancing pass must place more tracts into correct
districts after a coarse move.
The $200$-iteration boundary-swap rebalancer has a fixed budget; for large
$k$, coarse projections are more likely to exhaust the budget without
achieving balance, leading to rejection.
A lower $\alpha$ reduces the coarse rejection rate.

\paragraph{County population heterogeneity.}
States with highly unequal county populations (one or two dominant counties)
have a coarser population structure.
County-level proposals in such states tend to create imbalanced projections
and require more rebalancing iterations.
Again, a lower $\alpha$ is preferable.

\subsection{The Cost of Manual Tuning}

Manually tuning $\alpha$ across 50 states requires running short pilot chains
for each state, measuring the coarse acceptance rate, and adjusting $\alpha$
iteratively.
This is a significant practical burden:
\begin{itemize}
  \item Each pilot run takes minutes to hours for large states.
  \item The optimal $\alpha$ may shift as the chain evolves (non-stationarity).
  \item Retuning is required whenever the step count or district count changes.
\end{itemize}
A self-tuning procedure would eliminate this burden entirely.

\subsection{Robbins-Monro Stochastic Approximation}

Robbins and Monro~\citeyear{robbins1951} introduced a general method for
finding the root of an unknown function $h(\alpha) = 0$ using noisy
observations of $h$.
Applied here, $h(\alpha) = \mathrm{E}[\text{coarse accept rate}(\alpha)] - \text{target}$:
the expected coarse acceptance rate is (approximately) monotone in $\alpha$
over the operating range, so the root of $h$ is the $\alpha^*$ that achieves
exactly the target acceptance rate.

The Robbins-Monro update rule is:
\[
  \alpha_{t+1} = \alpha_t + \gamma_t \bigl(\hat{h}_t(\alpha_t) - 0\bigr),
\]
where $\hat{h}_t$ is the observed excess acceptance rate in the current window
and $\gamma_t = \gamma_0 / \sqrt{t}$ is a decaying step size.
This step-size schedule satisfies the classical conditions:
$\sum_{t=1}^\infty \gamma_t = \infty$ (sufficient exploratory power) and
$\sum_{t=1}^\infty \gamma_t^2 < \infty$ (convergence of variance to zero).
Under the monotonicity assumption, the iterate converges to $\alpha^*$
with probability one.

\subsection{Main Contributions}

This paper makes the following contributions:

\begin{enumerate}
  \item \textbf{Algorithm (\ams{}).}
    We define \ams{} (Section~\ref{sec:algorithm}), which wraps \ms{} with
    a Robbins-Monro adaptation loop that updates $\alpha$ every
    $\Delta = 50$ steps.
    Three design choices are analysed: the $\gamma_0/\sqrt{t}$ step-size
    schedule, the $\Delta = 50$ adaptation interval, and the $[0.05, 0.95]$
    clip to prevent degenerate fixed-point collapse.

  \item \textbf{Convergence analysis (Section~\ref{sec:convergence}).}
    We state a formal convergence theorem (under the monotonicity assumption),
    give a proof sketch citing Robbins-Monro 1951, and exhibit empirical
    alpha traces for NC, TX, and IA showing that convergence occurs within a
    few hundred steps.

  \item \textbf{Empirical comparison (Section~\ref{sec:comparison}).}
    We compare \ams{} to fixed-$\alpha$ \ms{} and to single-scale \recom{}
    on five states, measuring autocorrelation of the edge-cut objective
    $\EC$ at lag 100.
    \ams{} matches or improves on fixed-$\alpha$ \ms{} with no manual tuning.

  \item \textbf{Audit trail.}
    The $\alpha$ trace, $\gamma_0$, initial $\alpha$, and final $\alpha$ are
    recorded in the run manifest, enabling full reproducibility and independent
    verification of the convergence path.
\end{enumerate}

\subsection{Paper Structure}

Section~\ref{sec:background} reviews Multi-scale MCMC, Robbins-Monro
approximation, and adaptive MCMC theory.
Section~\ref{sec:algorithm} gives the \ams{} algorithm with pseudocode and
design rationale.
Section~\ref{sec:convergence} analyses convergence.
Section~\ref{sec:comparison} presents the empirical comparison.
Section~\ref{sec:conclusion} concludes.

\subsection{Relation to Prior Work}

\ams{} is a Search-layer algorithm, sitting at the same layer as \ms{} (G.11)
in the three-layer compositor.
It differs from the SeedCompositor strategies (\cs{}, \be{}) which change
\emph{how many} seeds are explored, not how each chain is run.
The adaptive MCMC literature~\citep{haario2001, andrieu2008} studies
adaptation of the \emph{proposal kernel} (e.g., covariance of a Gaussian
proposal); \ams{} adapts the \emph{mixing ratio} between two fixed kernels,
which is a structurally simpler problem with more direct convergence arguments.
